\documentclass[11pt,a4paper]{article}
\usepackage[utf8]{inputenc}
\usepackage{graphicx}
\usepackage[left=0.5cm,top=3cm,right=2.5cm,bottom=3cm,bindingoffset=0.5cm]{geometry}
\usepackage{macros/AEDLogica, macros/AEDEspecificacion, macros/AEDTADs}
\usepackage{caratula}


\titulo{Trabajo práctico}
\subtitulo{Especificacion TAD: Sistema de Fidelizaci´on AirMiles}

\fecha{\today}

\materia{Algoritmos y Estructuras de Datos}
\grupo{Grupo 1} 

\integrante{Garcia Ocio, Santiago}{1210/25}{santiagogarciaocio@gmail.com}
\integrante{Seda, Manuel} {631/18}{manu.seda@hotmail.com}
% \integrante{Apellido, Nombre2}{002/01}{email2@dominio.com} %


% Declaramos donde van a estar las figuras
% No es obligatorio, pero suele ser comodo
\graphicspath{{../static/}}

% Asi pueden escribir nuevos comandos. 
% Este por ejemplo asegura q los nombres 
% que figuren con una tipografia diferenciada  
\newcommand{\Tipo}[1]{\mathsf{#1}} 
% la sintaxis es \newcommand{\nombreDeLaMacro}[cantidadDeParametros]{Lo que va ser remplazado por el macro} 
\newcommand{\norm}[1]{\vert #1\vert}



\begin{document}

\maketitle

\section{Renombres}

\begin{itemize}
    \item[]NombrePromocion ES String
    \item[]PromocionActiva ES $\tupla{Factor, tieneTope, montoTope, millasAcumuladas}$
    \item[]PromocionesActivas ES $\seq{PromocionActiva}$
    \item[]Factor ES $\mathbb{R}$ 
    \item[]TieneTope ES Bool
    \item[]MontoTope ES $\mathbb{Z}$    
    \item[]Usuario ES $\mathbb{Z}$  
    \item[]HistorialVuelo ES $\seq{Vuelo}$
    \item[]HistorialTransaccion ES $\seq{Transaccion}$
    \item[]Vuelo ES $\tupla{Distancia,Millas,NombrePromocion,Factor}$
    \item[]Distancia ES $\mathbb{R}$
    \item[]Transaccion ES $\tupla{Accion X Millas X Comentario}$
    \item[]enum Accion ES \{\ "Transferir","Recibir","Canjear","Acumular","Referidor","Referido" \}\ 
    \item[]Comentario es String 
    \item[]Millas ES $\mathbb{R}$     
    \item[]enum Categoria ES \{\ "Bronce","Plata","Oro","Platino"\}\
    
\end{itemize}


\begin{tad}{AirMiles}


\obs{historialesVuelos}{\Diccionario{Usuario}{HistorialVuelo} }
\\
\obs{historialesTransacciones}{\Diccionario{Usuario}{HistorialTransaccion} }
\\
\obs{promociones}{\Diccionario{NombrePromocion}{PromocionActiva}}
\\


\begin{proc}{crearSistema}
{}{
\Tipo{AirMiles}
}
    \requiere{True}
    \aseguraLargo{
        res.promociones =  <"NoPromo":(1,False,0)>  \land \\
        res.historialesVuelos = <> \land \\
        res.historialesTransacciones = <> 
        }
\end{proc}

\begin{proc}{registrarUsuario}
{
\Inout am:AirMiles, 
\In usuario:Usuario
}
{}
    \requiere{am=AM_0   \land  usuario \notin AM_0.historialesVuelos }
    %no se si pedir tmb que no este en historialesTransacciones. Deberia bastar con chequear solo uno xq tienen las mismas claves los dos diccs
    \aseguraLargo{
        am.historialesVuelos = añadirUsuario(AM_0.historialesVuelos,usuario)\ \land \\
        am.historialesTransacciones = añadirUsuario(AM_0.historialesTransacciones,usuario)\ \land 
        am.promociones = AM_0.promociones
        }
\end{proc}

\begin{proc}{registrarUsuarioConReferidor}
{
\Inout am:AirMiles, 
\In usuarioReferido:Usuario, \\
\In usuarioReferidor: Usuario
}{} 
    \requiere{am=AM_0  \land usuarioReferidor \in AM_0.historialesTransacciones \land\\
    usuarioReferido \notin AM_0.historialesTransacciones} 
    \aseguraLargo{  
         am.historialesTransacciones[usuarioReferidor] =\\ agregarTransaccion(AM_0.historialesTransacciones[usuarioReferidor],\\("Refiridor",500,"Referido"))\\\land  \\
        igualesSalvoUsuariosIdaVuelta(AM_0.historialesTransacciones,  am.historialesTransacciones,\\\seq{usuarioReferidor, usuarioReferido}) \land \\  
        (usuarioReferido \in historialesTransacciones \yLuego 
        am.historialesTransacciones[usuarioReferido] =\\ \seq{("Referido", 500, "")})
         \land \\        
        am.historialesVuelos = añadirUsuario(AM_0.historialesVuelos,usuarioReferido) \land \\    
        am.promociones = AM_0.promociones 
                    }
\end{proc}

\begin{proc}{acumularMillas}
{
\Inout am:AirMiles,
\In usuario:Usuario,
\In distancia: $\mathbb{R}$,
\In promocion: NombrePromocion
}{}
    \requiere{am=AM_0\  \land \ usuario \in AM_0.historiales\  \land \ promocion \in AM_0.promociones\  \land \ distancia>0}
    \aseguraLargo{todasPromocionesIgualesSalvoUna(am, AM_0, promocion) \land \\
        igualesSalvoUsuarioIdaVuelta(AM_0.historialesVuelos, am.historialesVuelos, $\seq{usuario}$) \land \\
        igualesSalvoUsuarioIdaVuelta(AM_0.historialesTransacciones, am.historialesTransacciones, $\seq{usuario}$) \land \\
        (esPromocionConTope(am, promocion) \land promocionPasaTope(am, usuario, distancia, promocion) \implies am.historialesTransacciones[usuario] = agregarElemento(AM_0.historialesTransacciones[usuario],\\("Acumular",millasVuelo(am, usuario, distancia, "NoPromo"))) \land am.historialesVuelos[usuario] = agregarElemento(AM_0.historialesVuelos[usuario], (distancia, millasVuelo(am, usuario, distancia, "NoPromo"),\\ "NoPromo", 1))) \land 

        (esPromocionConTope(am, promocion) \land \lnot promocionPasaTope(am, usuario, distancia, promocion) \implies (am.promociones[promocion]_3 = am.promociones[promocion]_3 + millasVuelo(am, usuario, distancia, promocion)) \land am.historialesTransacciones[usuario] = agregarElemento(AM_0.historialesTransacciones[usuario],\\("Acumular",millasVuelo(am, usuario, distancia, promocion))) \land am.historialesVuelos[usuario] = \\agregarElemento(AM_0.historialesVuelos[usuario], (distancia, millasVuelo(am, usuario, distancia, promocion),\\ promocion, factorDePromocion(promocion)))) \land

        (\lnot esPromocionConTope(am, promocion) \implies am.historialesTransacciones[usuario] = \\agregarElemento(AM_0.historialesTransacciones[usuario],("Acumular",millasVuelo(am, usuario, distancia,\\ promocion))) \land am.historialesVuelos[usuario] = agregarElemento(AM_0.historialesVuelos[usuario],\\ (distancia, millasVuelo(am, usuario, distancia, promocion), promocion, factorDePromocion(promocion))))
        
    }
\end{proc}


\begin{proc}{canjearMillas}
{
\Inout am:AirMiles,
\In usuario:Usuario,
\In millas:Millas
}{}
    \requiere{am=AM_0\  \land \ usuario\in AM_0.historialesTransacciones[usuario]\  \land \ millas \leq millasDisponibles(AM_0.historialTransacciones[usuario])}
    %quizas es mas declarativo tener millasDisponibles(AM_0.historialTransacciones,usuario)
    %Con los corchete me hace 'ruido' tipo programacion y no tanto escribir logica
    \aseguraLargo{
        am.promociones = AM_0.promociones \land \\
        am.historialViajes = AM_0.historialViajes \land \\
        igualesSalvoUsuarioIdaVuelta(AM_0.historialTransaccions, am.historialTransaccions,\\ \seq{usuario}) \land \\
        am.historialTransacciones[usuario] = \\agregarTransaccion(AM_0.historialTransacciones[usuario], \tupla{"Canjear"XmillasX""})
                }
\end{proc}


\begin{proc}{transferirMillas}
{
\Inout am:AirMiles,
\In dador : Usuario,
\In receptor : Usuario,\\
\In millasTrans : Millas
}{}    
    \requiere{am=AM_0 \land dador \in AM_0.historialesTransacciones \land \\receptor \in AM_0.historialesTransacciones\ \land \\ millasDisponibles(AM_0.historialTransacciones[Dador]) \geq millasTrans}
    \aseguraLargo{am.promociones = AM_0.promociones \land \\
                am.historialVuelos = AM_0.historialVuelos \land \\
                igualesSalvoUsuarioIdaVuelta(AM_0.historialTransacciones, am.historialTransacciones,\\ \seq{dador,receptor}) \land \\
                am.historialTransacciones[dador] = agregarTransaccion(AM_0.historialTransacciones[dador],\\ \tupla{"Transferir"XmillasTransXreceptor}) \land \\
                am.historialTransacciones[receptor] = agregarTransaccion(AM_0.historialTransacciones[receptor],\\ \tupla{"Recibir"XmillasTransXdador})
                    }  
\end{proc}


\begin{proc}{consultarSaldo}
{
\In am:AirMiles,
\In usuario:Usuario
}{
\Tipo{Millas}
}
    \requiere{usuario \in am.historialTransacciones}
    \aseguraLargo{res = millasDisponibles(am.historialTransacciones[usuario])}
\end{proc}

\begin{proc}{consultarCategoria}{
\In am:AirMiles,
\In usuario:Usuario
}{
\Tipo{Categoria}
}
    \requiere{usuario \in am.historialesTransacciones}
    \aseguraLargo{(res = "Bronces" \iff totalDeMillas(am.historialesTransacciones[usuario]) < 10.000) \lor \\
            (res = "Plata" \iff (totalDeMillas(am.historialesTransacciones[usuario] \geq 10.000 \land totalDeMillas(am.historialesTransacciones[usuario]) < 50.000) \lor \\
            (res = "Oro" \iff (totalDeMillas(am.historialesTransacciones[usuario]) \geq 50.000 \land totalDeMillas(am.historialesTransacciones[usuario]) < 1.000.000) \lor \\
            (res = "Platino" \iff totalDeMillas(am.historialesTransacciones[usuario]) \geq 1.000.000) \\
            }
\end{proc}

\begin{proc}{rankingMillas}
{
\In am:AirMiles,
\In n:$\mathbb{Z}$
}{
\Tipo{\seq{Usuario}}
}
    \requiere{n > 0}
    \aseguraLargo{
        (n \geq |am.historialesTransacciones| \implicaLuego res = claves(am.historialesTransacciones)) \lor_l 
        ((\forall i \in \Z)(0 \geq i < n \implicaLuego rankingMillasDisponibles(am.historialesTransacciones,res[i]) \geq n)  \yLuego |res| = n )                      
    }
\end{proc}





\begin{proc}{resumenTransacciones}
{
\In am:AirMiles,
\In usuario:Usuario
}{
\Tipo{\tupla{MillasXMillasXMillas}}
}
    \requiere{usuario \in am.historialesTransacciones}
    \aseguraLargo{res_0 = acumuladoPorViajes(am.historialesTransacciones[usuario]) \land \\
                res_1=totalCanjeadas(am.historialesTransacciones[usuario]) \land \\
                res_2=totalTransferidas(am.historialesTransacciones[usuario])}
\end{proc}


\begin{proc}{iniciarPromocion}
{
\Inout am:AirMiles,
\In nombrePromocion : NombrePromocion,
\In factor : Factor
}{}
    \requiere{am=AM_0 \land factor \geq 1 \land nombrePromocion \notin AM_0.promociones}
    \aseguraLargo{am.historialesTransacciones = AM_0.historialesTransacciones \land \\
        am.historialesVuelos = AM_0.historialesVuelos \land \\
        am.promociones = añadirPromocionSinTope(AM_0.promociones,nombrePromocion,factor)
        }
\end{proc}


\begin{proc}{iniciarPromocionConTope}
{
\Inout am:AirMiles,
\In nombrePromocion : NombrePromocion, \\
\In factor : Factor,
\In montoTope : MontoTope
}{}
    \requiere{am=AM_0 \land factor >= 1 \land nombrePromocion \notin AM_0.promociones}
    \aseguraLargo{am.historialesTransacciones = AM_0.historialesTransacciones \land \\
        am.historialesVuelos = AM_0.historialesVuelos \land \\
        am.promociones = añadirPromocionConTope(AM_0.promociones,nombrePromocion,factor, montoTope)
    }
\end{proc}


\begin{proc}{finalizarPromocion}
{
\Inout am:AirMiles,
\In nombrePromocion : NombrePromocion
}{}
    \requiere{am=AM_0\ \land  \ bnombrePromocion \in AM_0.promociones\ \land \ nombrePromocion \neq "NoPromo"}
    \aseguraLargo{am.historialesTransacciones = AM_0.historialesTransacciones \land \\
        am.historialesVuelos = AM_0.historialesVuelos \land \\
        am.promociones = eliminarPromocion(AM_0.promociones,nombrePromocion)
    }
\end{proc}


\begin{proc}{elQueTieneMasMillas}
{
\In am:AirMiles
}{
\Tipo{Usuario}
}
    \requiere{True}
    %no se si esta bie npedir no vacio o si deberia tener que aceptar el vacio y devolver nada
    \aseguraLargo{
        res \in am.historialesTransacciones \land \\
        (\forall k \in \Z)(k \in am.historialesTransacciones\ \implicaLuego \ millasTotales(am.historialesTransacciones[k])\\ \leq millasTotales(am.historialesTransacciones[res]))
        }
\end{proc}


\begin{proc}{elMasCanjeador}
{
\In am:AirMiles
}{
\Tipo{Usuario}
}
    \requiere{True}
    %NO se aca si hay que pedir que haya al menos un elemento. No se si el vacio.
    \aseguraLargo{
        res \in am.historialesTransacciones \land \\
        (\forall k \in \Z)(k \in am.historialesTransacciones \implicaLuego totalTransferidas(am.historialesTransacciones[k]) \leq totalTransferidas(am.historialesTransacciones[res]))
    }
    
\end{proc}


\begin{proc}{paresTransferidoresExclusivos}
{
\In am:AirMiles,
}{
\Tipo{\seq{\tupla{UsuarioXUsuario}}}
}
    \requiere{True}
    \aseguraLargo{
        (|am.historialesTransacciones| < 2 \implicaLuego res = \seq{}) \land \\
        (sinRepetidosNiAnagramas(res) \land \\
        (\forall i,j \in Z)((i,j \in am.historialesTransacciones\ \yLuego i\neq j\ \land \\ sonParejaTransferidores(am.historieslTransacciones,i,j))\implicaLuego (\tupla{iXj} \in res) 
    }
\end{proc}

\pred{sonParejaTransferidores}{histoTrans: HistorialesTransacciones, usuario1:Usuario,\\ usuario2:Usuario}{\\
(\forall i \in Z)(0 \leq i < |histoTrans[usuario1]| \yLuego (histoTrans[usuario1][i]_0 = "Transferir"\ \lor \ histoTrans[usuario1][i]_0 = "Recibir")) \implicaLuego (histoTrans[usuario1][i]_2 = usuario2) \land \\
(\forall j \in Z)(0 \leq j < |histoTrans[usuario2]| \yLuego (histoTrans[usuario2][j]_0 = "Transferir" \lor histoTrans[j]_0 = "Recibir")) \implicaLuego (histoTrans[j]_2 = usuario1)
}
\\
\\
\pred {sinRepetidosNiAnagramas}{listaTuplaUsuarios:\seq{\tupla{UsuarioXUsuario}}}{\\
    (\forall i,j \in Z)(0 \leq i,j < |listaTuplaUsuarios|\ \yluego \ i\neq j\ \implicaLuego ((listaTuplaUsuarios[i]_0 \neq listaTuplaUsuarios[j]_0 \land listaTuplaUsuarios[i]_1 \neq listaTuplaUsuarios[j]_1) \land \\(listaTuplaUsuarios[i]_0 \neq listaTuplaUsuarios[j]_1 \land \ listaTuplaUsuarios[i]_1 \neq listaTuplaUsuarios[j]_0)))
}
\\
\\
\aux{millasDisponibles}{hist: HistorialTransaccion}{\R}{$\sum_{i=0}^{|hist-1|}$ fThenElse(esUso(hist[i]_0), -1*hist[i], hist[i])} 
\\
\\
\pred{esUso}{accion:Accion}{ accion \in  \seq{Canjear, Transferir} }
\\
\\
\aux{millasTotales}{hist: Historial}{\R}{\sum_{i=0}^{|hist-1|} IfThenElse(hist[i]_0 = "Transferir", -1*hist[i], hist[i])}
\\
\\
\aux{agregarTransaccion}{historialViejo: Historial,  transaccionNueva :Transaccion}{Historial}{historialViejo ++ <transaccionNueva>} 
\\
\\
\aux{añadirUsuario}{baseDeDatos:\Diccionario{Usuario}{\seq{T}}, usuario:Usuario}{\Diccionario{Usuario}{\seq{T}}}{setKey(baseDeDatos,usuario,\texttt{<>}) }
\\
\\
\aux{añadirPromocionSinTope}{baseDeDatos:\Diccionario{NombrePromocion}{\seq{T}}, promocion:NombrePromocion, factor:Factor}{\Diccionario{NombrePromocion}{\seq{T}}}{\\setKey(baseDeDatos,promocion,\tupla{Factor,False,0}) }
\\
\\
\aux{añadirPromocionConTope}{baseDeDatos:\Diccionario{NombrePromocion}{\seq{T}}, promocion:NombrePromocion, factor:Factor, tope : MontoTope}{\Diccionario{NombrePromocion}{\seq{T}}}{setKey(baseDeDatos,promocion,\tupla{Factor,True,tope}) }
\\
\\
\aux{eliminarPromocion}{baseDeDatos:\Diccionario{NombrePromocion}{\seq{T}}, promocion:NombrePromocion}{\Diccionario{NombrePromocion}{\seq{T}}}{delKey(baseDeDatos,promocion)}
\\
\\
\pred{igualesSalvoUsuarios}{dictControl : \Diccionario{Usuario}{\seq{T}} , dictAlterado : \Diccionario{Usuario}{\seq{T}} , usuariosModificados: \seq{usuario}}{(\forall k : \Z) (k \in dictControl \ \yLuego \ k \notin usuariosModificados))\implicaLuego (k \in dictAlterado \yLuego dictControl[k] = dictAlterado[k]) }
\\
\\
\pred{igualesSalvoUsuariosIdaVuelta}{dict1 : \Diccionario{Usuario}{\seq{T}} , dict2 : \Diccionario{Usuario}{\seq{T}} , usuariosModificados: \seq{usuario}}{
%Aca podriamos usar doble iguales SalvoUsuario, quizas es mas declarativo / les guste mas a los profes
(\forall k : \Z) (k \in dict1 \ \land \ k \notin usuariosModificado))\implicaLuego (k \in dict2 \yLuego dict1[k] = dict2[k]) \land \\
(\forall q : \Z) (k \in dict2 \ \land \ q \notin usuariosModificado))\implicaLuego (q \in dict1 \yLuego dict1[q] = dict2[q])
}
\\
\\
\aux{acumuladoPorViajes}{historial:historialesTransaccion}{Millas}{$\sum_{i=0}^{|historial|-1}$ ifThenElse(historial[i]_0 ="Acumular",historial[i]_1,0)}
\\
\\
\aux{totalCanjeadas}{historial:historialesTransaccion}{Millas}{$\sum_{i=0}^{|historial|-1}$ \\ifThenElse(historial[i]_0 ="Canjear",historial[i]_1,0)}
\\
\\
\aux{totalGastadas}{historial:historialesTransaccion}{Millas}{$\sum_{i=0}^{|historial|-1}$ \\ifThenElse(historial[i]_0 \in \seq{"Canjear","Transferir"},historial[i]_1,0)}
\\
\\
\aux{totalTransferidas}{historial:historialesTransaccion}{Millas}{$\sum_{i=0}^{|historial|-1}$ \\ifThenElse(historial[i]_0 \in \seq{"Transferir","Recibir"},historial[i]_1,0)}
\\
\\
\aux{rankingMillasDisponibles}{historial:HistorialTransacciones, usuario:Usuario}{\Z}{ |historial| + 1 - \\$\sum_{\forall k \in historial}$ ifThenElse(millasDisponibles(historial[k]) \leq millasDisponibles(historial[usuario]),1,0)
\\
\\
\pred{esPromocionConTope}{am : AirMiles, promocion : NombrePromocion}
{am.promociones[promocion]_1 = True}
\\
\\
\pred{promocionPasaTope}{am : AirMiles, usuario : Usuario, distancia : Distancia, promocion : NombrePromocion}{millasVuelo(am, usuario, distancia, promocion) + am.promociones[promocion]_3 > am.promociones[promocion]_2}
\\
\\
\aux{agregarVuelo}{am : AirMiles, usuario : Usuario, distancia : Distancia, promocion : NombrePromocion}{AirMiles} = Agregar
% no sé bien cómo escribir esto

}



\end{tad}


\end{document}
